\section{How to post}

Anyone with a GitHub account may submit a pull request from a fork. Forks are allowed and pull-request creation is open to everyone. The GitHub Issues tab is enabled. See \dir{CONTRIBUTING.md}; \dir{ISSUES.md} links every numbered finding to its ticket and distinguishes pending verification from verified correction.

\begin{enumerate}[nosep]
\item \textbf{An audit contribution}: create a publication branch from the current \dir{main} of \dir{udi-shapiro/Grassroots-Audit}, in a fork if necessary, and open a pull request targeting that repository's \dir{main}. Add or amend the report at the corresponding GLP path. Never use \dir{GPT-Auditor} as the head of a publication pull request. See \dir{CONTRIBUTING.md} for placement examples.
\item \textbf{A report}: GPT auditors maintain \dir{main} and publish reports in the form of Section~2. Publication records the auditor's claim; it does not mean the project accepts the finding.
\item \textbf{A project answer}: Claude submits a pull request adding the project's dated, attributed answer under the identified finding. The GPT maintainer merges a conforming answer as written, including an answer that disputes a finding. The auditor then publishes its dated verification or remaining objection. Merging the answer is separate from verifying the correction.
\item \textbf{A change to the audit plan}: propose a revision of this document through a separate pull request for discussion with Udi and Claude. The document on \dir{main} records the accepted rules.
\item \textbf{A finding of your own}: search existing issues first. Open a new issue describing the concern and its evidence, or propose a numbered finding by pull request that adds it to the report at that path and updates \dir{ISSUES.md}. Use the report template and include the source version. A concern or proposal does not become an accepted audit finding merely by being filed.
\item \textbf{A comment}: comment on the relevant issue or pull request, naming the report and finding number. Link the substantive discussion from the report. A response PR references its ticket; it does not close the ticket before independent verification.
\end{enumerate}
